% !PRINTED 392
\weberpage{392}
%<<CONT
$\omega_n$ a quadratic equation for $\omega$, which can be put in the form:
\[
\begin{vmatrix} P_n - \omega Q_n, & P_{n-1} - \omega Q_{n-1} \\ P_m - \omega Q_m, & P_{m-1} - \omega Q_{m-1} \end{vmatrix} = 0.
\]

\begin{center}
\S.\ 126.

Examples.
\end{center}

We shall here carry out some examples of the determination of the reduced
numbers and periods, in order to show the application
of the method; the examples can of course be multiplied
quite at pleasure.

1. $D = 29$.

We first determine, according to \S.\,125, 4., all the
reduced numbers belonging to $D$. The greatest integer $\lambda$
contained in $\sqrt{D}$ is here $= 5$; hence $b$ can have only the values:
\[
b = 1, 3, 5,
\]
whence
\[
ac = \frac{D-b^2}{4} = 7, 5, 1.
\]

The limits within which $a$ and $c$ must lie are, by
\S.\,125, (7), for the three values of $b$:
\[
3, 3; \quad 2, 4; \quad 1, 5.
\]
One therefore obtains only the single reduced number:
\[
\{1, 5, 1\}.
\]

The period consists of a single term, and the
continued fraction for $\omega$ is:
\[
\omega = \frac{5+\sqrt{29}}{2} = (5, 5, 5 \ldots)
\]

2. $D = 116 = 4.29$.

Here $\lambda = 10$ and $b$ has one of the values
\[
b = 2, 4, 6, 8, 10,
\]
hence
\[
ac = 28, 25, 20, 13, 4;
\]
the limits for $a$ and $c$ are
\[
5, 6; \quad 4, 7; \quad 3, 8; \quad 2, 9; \quad 1, 10,
\]
hence, since the cases in which $a$, $b$, $c$ have the common factor
%CONT>>
